\chapter*{Appendices}
\addcontentsline{toc}{chapter}{Appendices}

\section*{Appendix 1: Channel Layout of the Feature Vector}

Each 25\,Hz sample is a 45-channel vector in finger-major order. For each finger
in the order thumb, index, middle, ring, pinky, the nine fields are: pitch (P),
roll (R), yaw (Y), angular rate about $x$/$y$/$z$ ($g_x,g_y,g_z$), and linear
acceleration along $x$/$y$/$z$ ($a_x,a_y,a_z$). This ordering is identical in
the data-collection script, the master firmware payload, the normalisation
statistics, and the mobile application's parser.

\section*{Appendix 2: Vocabulary}

\textbf{Letters:} the 26 letters A--Z. Most are encoded as short, distinctive
finger motions; J and Z are the inherently dynamic ASL letters.

\textbf{Words (19 one-handed signs):} DRINK, EAT, FINE, HELLO, HOWAREYOU,
ILOVEYOU, ME, NAME, NO, PLEASE, SORRY, STOP, THANKYOU, WATER, WHAT, WHERE, YES,
YOU, YOUR. The compact multi-word tokens (e.g.\ HOWAREYOU, ILOVEYOU) are
expanded to spaced phrases (``HOW ARE YOU'', ``I LOVE YOU'') for display and
speech in the application.

\section*{Appendix 3: Key System Parameters}

\begin{tabular}{ll}
\toprule
Parameter & Value \\
\midrule
Fingers / channels        & 5 / 45 \\
Sampling rate             & 25\,Hz \\
Capture window            & 3\,s (75 time steps) \\
ESP-NOW channel           & 11 \\
BLE MTU (negotiated)      & 256 bytes \\
BLE payload               & 180 bytes (45 float32) \\
Master CSV baud rate      & 921600 \\
Model quantisation        & int8 (input and output) \\
\bottomrule
\end{tabular}

\section*{Appendix 4: Source Code}

The firmware, data-processing scripts, and mobile application are organised in
the project repository under \texttt{firmware/}, \texttt{scripts/}, and
\texttt{flutter\_app/} respectively.
